\section{What Studio does and where it stops}
Nimbus Studio is the desktop workspace for setting up atmospheric experiments,
launching the Nimbus executable, inspecting saved fields, comparing datasets and
publishing figures. It gives one experiment a common place for its configuration,
run status, data sources, plots and analysis recipes. The intent is to make careful
experimentation practical on ordinary computers without requiring each researcher
to rewrite a model or build a separate visualization application.

The atmospheric solver and Studio are independent programs. The solver advances
the equations; Studio displays products the solver has published. A camera move,
color change or Python analysis does not advance or modify the simulated state.
Closing a plot does not delete the simulation's output. Conversely, a pleasing
render is not evidence that its source experiment satisfies a conservation,
convergence or scientific-validation requirement.

Nimbus was independently developed from the beginning. CM1 appears in Studio as a
read-only dataset and configuration-interchange source, not as the engine behind
Nimbus. Similar variable names or familiar controls do not imply identical
equations, initial conditions, closures or numerics. Use the
\href{MATHEMATICS.pdf}{mathematics guide} for those assumptions and the
\href{ARCHITECTURE.pdf}{architecture guide} for solver ownership and execution.

This guide follows the ordinary application paths. It describes implemented
features and their admission checks, including intentionally unavailable paths.
It does not certify every machine, renderer, field combination or scientific
interpretation. The terms Ready and Fully Qualified retain the repository's
meaning: bounded tests, completed long campaigns and independent scientific
review are distinct evidence. A successful image is only one part of that chain.

\begin{figure}[htbp]
\centering
\begin{tikzpicture}[
  box/.style={draw=NimbusBlue,rounded corners,fill=NimbusLight,align=center,
    text width=4.8cm,minimum height=1.0cm,font=\small},
  flow/.style={-{Latex[length=2mm]},thick,draw=NimbusBlue}]
\node[box] (config) {Configuration and executable catalog\\Select, validate, launch};
\node[box,right=1.0cm of config] (solver) {Nimbus solver process\\Advance physical state};
\node[box,below=0.8cm of solver] (history) {Published history\\Completed frames\\Geometry and units};
\node[box,left=1.0cm of history] (reader) {Rust readers and bounded IPC\\Validate, reserve and decode};
\node[box,below=0.8cm of reader] (worker) {Frontend workers\\Samples and geometry\\Explicit plot ownership};
\node[box,right=1.0cm of worker] (screen) {Plotly or VTK and the graphics API\\Render, commit, capture};
\draw[flow] (config) -- (solver);
\draw[flow] (solver) -- (history);
\draw[flow] (history) -- (reader);
\draw[flow] (reader) -- (worker);
\draw[flow] (worker) -- (screen);
\end{tikzpicture}
\caption{The solver publishes data across a process boundary. Studio renders
those data through a separate ownership chain. Imported histories enter at the
reader boundary; changing the view does not feed back into the solver.}
\end{figure}


\subsection{The first launch has no storm to display}
A fresh checkout and fresh Studio profile contain no pre-run simulations. The
Runs view may be empty; the Plots workspace offers \textbf{Open data}. These are
normal states. Studio can open, show configuration controls and report its
preferences without running the solver or downloading scientific runtime tables.
The \href{BUILDING.md}{build guide} describes a build-and-open-only check for
Windows, Linux and macOS. This is the appropriate first check on a new machine.

There are then two ordinary ways to obtain a dataset. Open an existing Nimbus
run or supported CM1 history through Open data, or configure and launch a Nimbus
experiment. A run directory is the useful unit of Nimbus data: its metadata and
completion records give meaning to the field files. Selecting an arbitrary array
without that context cannot reconstruct its coordinate system, units, valid
domain or simulation time.

The native desktop application provides filesystem access, process launch and
native readers. The development browser frontend can exercise rendering and
synthetic harnesses, but its mocked native operations are not a replacement for
the installed application. Browser rendering evidence and native application
evidence should be reported separately.

\subsection{The main working areas}
\textbf{Configure} holds the experiment draft and its launch controls.
\textbf{Runs} exposes launched or discovered histories, status and logs.
\textbf{Plots} opens a data source and builds a viewing workspace. Settings
contains local paths, appearance, loading performance and remote connections.
Projects and investigation workspaces retain related source and view choices so
the user can return to an analysis without treating one transient screen as the
whole experiment.

Preferences such as text size, light or dark appearance and which side holds the
plot affect the desktop workspace. They are not solver configuration. Similarly,
the Loading workers preference controls Studio's field-loading concurrency,
whereas the solver's threading and compute-backend settings belong to the run.
Keep these two kinds of tuning separate when diagnosing performance.

Saved preferences and recent-dataset entries are conveniences, not archives of
the underlying data. Moving or deleting a source can make a saved project
unresolvable. Preserve the source history, configuration and any external inputs
alongside a project when reproducibility matters. Exported investigation and
research packages add explicit recipes and provenance; they do not silently turn
every linked source into an embedded simulation archive.

\section{Configuration, launch and data discovery}
\subsection{A draft must agree with the executable}
The configuration editor combines a structured form with YAML. Public templates
provide starting experiments; they are not certificates for a storm, a resolution
or an arbitrary set of physics choices. The
\href{NAMELIST.md}{complete namelist reference} inventories strict-v1 settings,
omission defaults, units, aliases and restrictions, including options not exposed
as form controls and legacy-reader compatibility.

Studio obtains a scheme catalog from the selected Nimbus executable. The catalog
describes its selectors and admission constraints. A UI label alone does not
establish that the executable implements the requested combination. Validation
and preflight check the actual draft, target, executable identity, required
inputs and configuration dependencies before launch. A configuration can contain
individually recognizable schemes while their coupling is still rejected.

Scientific tables are separate runtime inputs. Thompson and RRTMG requirements
are described in the build guide; opening Studio does not imply that these
tables have been installed. Provisioning a canonical table set and checking its
identity prevents compiler-dependent reconstruction from silently substituting
different input bytes. A missing input should be repaired at its reported path,
not bypassed by changing an unrelated numerical tolerance.

The binary-path preference selects an executable when automatic discovery is
not the intended build. The runs-directory preference selects the library root
for newly launched histories. Inspect these settings when several checkouts or
installed versions coexist: the source tree being edited and the executable
being launched are not necessarily the same build.

\subsection{Run status and continuation}
Launching creates a separate solver process. Studio collects status and log
information; run progress and saved-frame availability are different signals.
A process may be advancing while no new output frame is complete. Conversely,
a finished process can leave useful completed frames even when a later operation
failed. Read the reported outcome and logs rather than interpreting the presence
of a directory as success.

Checkpoint inspection belongs to continuation, not to plotting. A history frame
contains display products; a valid continuation must retain all state required
by the selected runtime. Do not assume that being able to plot a time means it
is a restart endpoint. The executable's checkpoint and configuration checks
remain authoritative when preparing a continuation.

\subsection{Nimbus histories and imported histories}
Nimbus output readers support the maintained NPY, ZFP and netCDF history paths.
NPY remains supported while the required netCDF replacement scope is being
completed and qualified. ZFP is a storage codec, not a new coordinate system or
a guarantee of losslessness: retain the output's compression provenance when
comparing values. Opening a compressed field should not erase how its values
were encoded.

A published output hierarchy can provide checksummed levels of detail and
bounded native reads. Without such a hierarchy, level zero still has its native
read contract; Studio must not invent a published coarse product merely because
it can produce a smaller display. Source selection records whether a plot uses
a published level, a bounded source window, a reduction or display reconstruction.

CM1 import can inspect history files, a case directory, namelist information and
an input sounding. Its compatibility and decision views distinguish direct
mechanical mappings from unresolved scientific decisions. Translation can produce
a partial draft that still needs interpretation. Nimbus does not install or run
CM1, and importing CM1 data does not make Nimbus a CM1 fork.

An adapter records source kind, locator, geometry and availability. Rank matters:
a two-dimensional column product cannot enable a three-dimensional renderer just
because its field name looks atmospheric. Vector views also require physically
matched components. The plot picker therefore uses dataset capabilities rather
than presenting every plot as if every file could satisfy it.

\subsection{Health and alignment are part of opening data}
Dataset health checks expose missing fields, shape or coordinate disagreements,
time-axis problems and incomplete products. Treat a missing value as missing,
not as calm wind, dry air or zero precipitation. A finite array value is not by
itself proof that the location lies inside the physical domain: halos and boundary
closure planes are excluded according to the saved geometry.

For comparisons, inspect the alignment report before interpreting a difference.
Time, physical coordinates, units, component meaning and valid-domain support
must match the requested operation. A common shape is insufficient: two arrays
can have equal dimensions but different origins, spacings, model frames or units.
When an operation requires aligned grids, a rejected mismatch is useful feedback,
not a reason to strip the metadata.

\subsection{A practical first inspection}
Open a completed history and start with one scalar Slice. Check its field name,
units, saved time, axis labels, slice location and source disclosure. Move to a
second completed frame and confirm that both the values and time label change.
Then add a vertical profile or cross section through an identifiable feature.
Only after those coordinate checks should a volume or composite storm scene be
used to explain the feature's three-dimensional appearance.

This sequence makes a dense 3-D view easier to interpret. An opaque volume can
hide an interior maximum; a threshold surface omits everything below its chosen
level; a derived diagnostic can require fields that were never exported. A simple
slice provides a direct way to distinguish a data problem from a camera,
transfer-function or feature-selection problem.
